\section{Results}
\label{sec:results}

\begin{figure}[t]
\centering
\paperfig{geometry.pdf}{1.0}
\caption{\textbf{Where the wall goes.} Left: squared latent distances in the grid JEPA (model seed $0$) over $64$ pre-reveal contexts; the wall's amplitude (latent of the true future with the wall against without it)
lies below the predictor's one-step residual at every step, and the rollouts drift far from the true
future. Middle: the flow model's latent on the axis between the two gap futures ($40$ contexts in which both branches survive); a sample commits at once, the mean of eight starts half-way and approaches a branch in three steps. Right: termination head fitted on the frozen JEPA latent, AUC of each death cause against alive transitions.}
\label{fig:geometry}
\end{figure}

\subsection{The wall is decodable but lost in prediction}
\label{sec:results-wall}

The wall's signal in the JEPA latent is smaller than the predictor's error
(Figure~\ref{fig:geometry}, left). From $64$ contexts in which the wall is solid at $x\in\{42,45\}$ the
true simulator renders three futures under the recorded actions: the recorded one, its counterfactual
with the other gap, and the same trajectory with the wall removed. The encoder separates
the true future from the wall-free one by a squared distance of $0.035$ at step one and $0.12$ at
step eight, in a $512$-dimensional latent of about unit variance per coordinate. The predictor's own teacher-forced one-step residual is
$0.145$ at step one and $0.25$ at step eight. Write $d$ for the difference between the latents of the
true future with and without the wall, and the wall coordinate of a prediction for its projection on
$d$, $0$ at the wall-free latent and $1$ at the true one: the teacher-forced one-step prediction has
mean wall coordinate $0.39$ at step one and $0.21$ at step eight. The deterministic rollout is as close to the
wall-free future as to the true one at every step, and its decoded frames contain no wall at any step,
while the decoded latent of the true future contains the full wall (intensity $1.00$ on the wall
pixels; Table~\ref{tab:latent}). Averaging cannot explain this: the wall is dropped before the gap is
drawn, and the two possible futures share the wall. These are model seed $0$; seeds $1$ and $2$ repeat
it with wall amplitude $0.03$ and $0.09$ against a residual of $0.13$ and $0.16$ at step one, wall
coordinate $0.45$ and $0.66$, termination heads at AUC $0.65$ and $0.76$ on wall deaths and agents at
$0.17$ and $0.33$ walls; seed $2$ keeps $0.68$ of the wall's decoded intensity at step one and at most $0.02$ after (Table~\ref{tab:seeddiag}).

The flow model keeps the wall in both modes. In its latent the same distance is $344$ to $554$; on
the $40$ contexts in which both branches survive the recorded actions, a sample sits at
$\lvert\text{gap coordinate}\rvert = 0.89$ at step one and $0.99$--$1.00$ after, the mean of eight at
$0.28$, $0.47$, $0.83$ and $0.98$ (Figure~\ref{fig:geometry}, middle; Table~\ref{tab:latent}). Its decoded frames show one gap or the other, never a half-open wall: a latent average need not decode to a pixel average.

The heads fitted on the latent see border deaths and barely see wall deaths
(Figure~\ref{fig:geometry}, right; Table~\ref{tab:heads}). Fitted on $80{,}000$ real latents with the
death labels, the termination head separates the $159$ border deaths from alive transitions at AUC
$0.963$ and the $610$ wall deaths at $0.649$ (model seed $0$; $0.88$--$0.96$ and $0.65$--$0.76$ over
seeds); the reward head predicts $-0.89$ on border deaths and $-0.05$ on wall deaths, against $-1$ at every death. These are the heads the planners and in-model agents use, not a bound on any read-out.

\subsection{Nothing trained inside the JEPA learns to survive}
\label{sec:results-ppo}

\begin{table}[t]
\centering
\caption{\textbf{Gate controllers, graded in the real environment.} Walls per episode and survival to the cap, mean $\pm$ s.d.\ over runs; ``discovered'' counts runs with final survival above $50\%$. Simulator search uses the true reward, no terminal value; real-gate PPO uses $1.5$M steps.}
\label{tab:paired}
\inputtable{paired.tex}
\end{table}

\begin{figure}[t]
\centering
\paperfig{discovery.pdf}{1.0}
\caption{\textbf{Survival in the real gate over policy training}, one line per run; discovery is final survival above $50\%$ (dashed line). Inside the grid JEPA no run leaves the floor. Inside the flow model and on true pixels with the convolutional torso some runs discover the retreat-and-wait strategy (dotted: the MLP torso on pixels, which does not in three runs). On the frozen JEPA latent no run discovers it; on the latent decoded back to pixels with the convolutional torso (dotted), $2$ of $9$ do. Walls passed, the quantity of the main comparison, are in Figure~\ref{fig:hero}C and Table~\ref{tab:paired}.}
\label{fig:discovery}
\end{figure}

Table~\ref{tab:paired} and Figure~\ref{fig:discovery} give the control result. PPO inside the grid
JEPA ends at or below $3.1\%$ survival ($8$ of $256$ episodes) at the final checkpoint in every one of its $27$ runs: $0.27\pm0.22$ walls per episode with
the original predictor, $0.26\pm0.30$ with the residual predictor as the mean of eight and
$0.26\pm0.32$ as one sample (Table~\ref{tab:seeds} lists every cell). The best cells, $0.9$ walls at $3\%$ survival, sit between horizon-5 and horizon-8 simulator search, and no run leaves the floor of Figure~\ref{fig:discovery}. Inside the flow model the eight runs without discovery pass $0.72$ to $1.18$ walls per episode, and one of the six sampled runs discovers the
retreat-and-wait strategy and reaches $98\%$ survival and $3.95$ walls; PPO in the real environment
discovers it in $2$ of $3$ runs ($91$, $5.5$ and $99\%$). Discovery is a lottery in both. Discovery counts are $1$ of $9$ for the flow model and $0$ of $27$ for the JEPA (Fisher $p=0.25$); all nine flow runs exceed the JEPA's mean of $0.27$ walls and six exceed its best run (Mann--Whitney $p<10^{-4}$). These run-level tests are conditional on the fitted models, and none of the $27$ JEPA policies discovers the strategy.

Two controls in the real gate, with the true dynamics and the true reward, show a deficit relative to training on real transitions. PPO on the frozen JEPA latent (Figure~\ref{fig:discovery}, right) ends at $0.88$--$0.99$ walls per episode and $2.3$--$5.5\%$ survival in nine runs over three encoders, $0.94$ against $0.27$ inside the JEPA, and never the strategy. The same MLP torso design passes $0.84$--$1.00$ walls on the flattened true frame in three runs, so the one-wall result is not specific to the JEPA latent. In the real gate the convolutional policy discovers the strategy in $2$ of $3$ pixel runs and, fed the frozen latent decoded back to pixels by the off-budget decoder, in $2$ of $9$; both decoded successes use encoder $1$ ($3.6$--$3.9$ walls, $79$--$96\%$; $0.80$--$1.39$ walls and $2$--$7\%$ in the other seven), so the frozen latent supports discovery on at least one encoder. None of $42$ transformed-latent runs discovers it; they pass $0.50$--$1.01$ walls and test wall-signal gains of $4$ and $16$, four-fold scaling of probe, wall-difference and matched random subspaces, and decoded frames for the MLP torso (Appendix~\ref{app:latentctl}). The comparison that carries our claim is within one policy class, $0.94$ walls on the frozen latent in the real gate against $0.27$ inside the JEPA; real-gate training also replaces the fitted heads by the true reward and changes the collection protocol (Appendix~\ref{app:hparams}), so it measures the imagination pipeline, predictor and heads together, and does not separate transition error from head error.

Two references bound what the planners can say. Simulator search with the wall known, the gap hidden
until the reveal and no terminal value passes $0.69$ walls at horizon $5$ and $1.03$ at horizon $8$;
revealing the gap to horizon-5 search moves it only to $0.70$: the horizon binds, not the gap.
LeWM's goal-image and heads planners on the JEPA pass $0.04$ to $0.16$ walls
(Appendix~\ref{app:round1tables}); in a horizon-$3$ audit the heads-MPC's selected return, $-0.46$, is a maximum over $64$ noisy plans with mean $-0.84$ (Table~\ref{tab:heads}).

\subsection{Sampling the fitted residual predictor does not rescue control}
\label{sec:results-mean}

The comparison is inside one model. In the grid JEPA the residual predictor has lower held-out one-step error than the original ($1.3\times10^{-4}$ per coordinate as the mean of eight, $2.1\times10^{-4}$ as one sample, against $3.1\times10^{-4}$) and is a distribution rather than a point; neither moves the agent: $0.26$ and $0.26$ walls against $0.27$, the same $0$--$3\%$ survival (Table~\ref{tab:paired}).
No distributional predictor is ruled out; this one does not help. In the flow model the mean of eight passes $0.88\pm0.14$ walls in three runs, against $0.96\pm0.16$ for the five sampled runs without discovery and $3.95$ for the one with it; $0$ against $1$ discoveries in three against six runs says nothing
about the mode. Proposition~\ref{prop:regret} remains true and Figure~\ref{fig:geometry} (middle) shows the half-way latent it warns about, but the averaged transition approaches a branch within three steps (gap coordinate $0.83$) and keeps the wall; inside the JEPA it is gone in either mode. These tests neither identify the mean as the cause nor rule out failures caused by averaging.

\subsection{Budget, regularizer, and the cube worlds}
\label{sec:results-rest}

\paragraph{Training budget.} Training the grid JEPA for twice and four times the steps ($17{,}000$
and $34{,}000$) lowers the prediction loss from $0.0006$ to $0.0003$ and leaves the heads' held-out reward correlation at $0.31$--$0.33$ and PPO inside the model at $0$--$3\%$ survival in all three predictor variants (model seed $0$, PPO seed $0$; Table~\ref{tab:budget}).

\paragraph{Regularizer dose.} On the cube world with a random texture fixed within each episode,
where the class-token JEPA behaves like the texture optimum of Proposition~\ref{prop:representation},
sweeping the SIGReg weight over $0$, $0.009$, $0.09$ and $0.9$ and the latent width over $32$ and $192$ (five arms) does not recover the cube position: the held-out reward correlation of the fitted heads stays between $-0.01$ and $0.00$ and the decoder's error at $0.0028$--$0.0029$ in every regularized arm; at weight $0$ the latent collapses (decoder error $0.055$; Table~\ref{tab:sigreg}). The weight sets how Gaussian the latent is, not what it is a function of.

\paragraph{Cube worlds and a state-blind caveat.} The six cube worlds (Appendix~\ref{app:cubes}) show
the two representation failures the theory predicts: with a per-episode random background the class-token JEPA reaches near-zero prediction loss by encoding the background and its reward-driven controllers fall to $0$--$2\%$; with a hidden wind the grid JEPA stays in the noise basin. A policy
that alternates left and right without looking survives the still cube worlds in $100\%$ of episodes (Table~\ref{tab:edge}); still-world survival is not evidence.

\paragraph{Craftax, exploratory.} A grid JEPA trained on Craftax-Classic pixels reaches six-action goal
positions in $94.5\%$ of trials against $18.8\%$ for random plans, while its reward head's correlation with the true reward falls from $0.92$ on training transitions to $0.35$ held out and reward-MPC returns $0.47$ over $100$ steps against $1.45$ for random actions (Appendix~\ref{app:craftax}): what it keeps steers the agent and does not say what is worth reaching.
